\SourceLeaf{523}\par\medskip\Needspace{4\baselineskip}\pdfbookmark[1]{372.}{sectio-octava-art-372}\centerline{\textbf{372.}}\nobreak\smallskip
Si $n$ est numerus primus (impar), residua $m$ consistent ex cifra, pro qua $M=1$, et ex $\frac12(n-1)$ aliis numeris, pro quibus $M=2$. Designando haec residua (excluso residuo $0$) indefinite per $\mu$, erit progressio nostra $=1+2\Sigma r^\mu$. Porro si per $\nu$ designantur indefinite omnes reliqui numeri infra $n$, quorum multitudo quoque erit $\frac12(n-1)$ et qui omnia non-residua quadratica ipsius $n$ infra $n$ complectentur, manifesto erit
\[
1+\Sigma r^\mu+\Sigma r^\nu=1+r+rr+r^3\dotsfour+r^{n-1}=\frac{1-r^n}{1-r}=0
\]
Quare ponendo summam progressionis nostrae sive $1+2\Sigma r^\mu=A$, erit $1+2\Sigma r^\nu=-A$, nec non $\Sigma r^\mu-\Sigma r^\nu=A$.
\par Per art. $356$ fit itaque $A=\pm\surd n$ vel $\pm\surd{-n}$, prout $n$ est $\equiv1$ vel $\equiv3\, (\mathrm{mod}.4)$. Sed signum radicis hinc nondum determinatur.
\par Si in progressione nostra, quam per $\Pi$ designabimus, pro $r$ substituitur alia similis radix aequationis $x^n-1=0$, puta $r'=r^\lambda$, supponamus inde prodire $\Pi'$.
\par\medskip\Needspace{4\baselineskip}\pdfbookmark[1]{373.}{sectio-octava-art-373}\centerline{\textbf{373.}}\nobreak\smallskip
Si $n$ est quadratum altiorve potestas numeri primi, puta $=p^\pi$, residua $m$ quaedam consistent e numeris per $p$ non divisibilibus, alia erunt divisibilia per $pp$ neque per altiorem potestatem ipsius $p$, alia per $p^4$ neque vero. per $p^5$ dividi poterunt et sic porro usque ad ea quae per $p^{\pi-2}$ neque vero per $p^{\pi-1}$ divisibilia sunt, sive per $p^{\pi-1}$ neque vero per $p^\pi$, prout $\pi$ par est sive impar; his denique accedit residuum $0$, quod est unicum per $p^\pi$ divisibile (conf. art. $102$). Iam designando per $\mu$ indefinite residua quadratica numeri $p$ infra $p$ cifra exclusa (quorum multitudo $=\frac12(p-1)$), illae diversae residuorum classes sequenti modo exhibebuntur. Prima, quae per $p$ non sunt divisibilia, repraesentantur per $\mu+kp$, ubi pro $k$ substituendi sunt omnes integri a $0$ usque ad $p^{\pi-1}-1$, ita ut omnium residuorum in hac forma contentorum multitudo sit $=\frac12(p-1)p^{\pi-1}$; pro his singulis fit $M=2$. Summa autem omnium terminorum in $\Pi$ his residuis respondentium erit
\[
\begin{aligned}
&=2\Sigma r^{\mu+kp}=2\Sigma r^\mu\pdot\Sigma r^{kp}\\
&=2\Sigma r^\mu\pdot\frac{r^{p^\pi}-1}{r^p-1}=0
\end{aligned}
\]
Secunda residuorum classis exhibetur per $\mu pp+kp^3$ ubi pro $k$ substituendi sunt omnes integri a $0$ usque ad $p^{\pi-3}-1$ ita ut omnium residuorum in hac
